
\chapter{Visa}\label{ch:visa}
В~общей четырёхсторонней модели
(см.~раздел~\ref{sec:ecosystem-four-party}) Visa задаёт собственную сетевую
архитектуру, свод правил, сервисы и~требования к~приёму карт.

\importantnote{%
  \textbf{Контекст для~российского читателя.} 5~марта 2022~года
  Visa приостановила работу с~российскими банками для~новых
  международных операций~\cite{visa-russia-suspension-2022}.
  Карты Visa, выпущенные российскими банками до~этого момента,
  продолжают приниматься внутри России: внутрироссийские
  авторизации и~клиринг по~таким картам идут через
  операционный платёжный и~клиринговый центр Национальной системы
  платёжных карт (ОПКЦ НСПК; см.~гл.~\ref{ch:mir}), минуя
  VisaNet~--- глобальную процессинговую сеть Visa, и~срок действия карт российские
  банки продлевают самостоятельно~\cite{cbr-russia-cards-normal-2022}.
  Международный приём этих карт за~пределами РФ прекращён.
  Глава описывает архитектуру и~правила Visa, на~которые российские эквайеры
  и~эмитенты ориентируются из-за наследия интеграций и~требований
  совместимости; российская специфика после 2022~года разбирается
  в~главе~\ref{ch:mir}.

}

\section{VisaNet: авторизационная сеть и~расчётный бэкенд}\label{sec:visa-visanet}

Авторизация и~клиринг разнесены по~времени и~характеру задачи
(см.~гл.~\ref{ch:tx-lifecycle}), поэтому Visa обслуживает их в~двух
независимых сетях.
Авторизационная сеть исторически называлась \textbf{BASE~I}. В~современных
документах Visa она именуется \textbf{V.I.P.~System} (Visa Integrated Payments)
и~включает не~только BASE~I (авторизации из~двух сообщений,
\textit{dual-message}), но и~\textbf{SMS (Single Message System)} для~сценариев
с~одним сообщением (\textit{single-message}), исторически~--- для~дебетовых
и~предоплаченных программ~\cite{visa-vip-system-overview-2005}. Сеть
отложенного клиринга и~расчёта по-прежнему называется \textbf{BASE~II};
денежный расчёт между участниками выполняет отдельная подсистема
\textbf{VisaNet Settlement Service (VSS)}, логически отделённая от~клирингового
конвейера BASE~II~\cite{visa-core-rules,visa-vip-system-overview-2005}.
\highlight{V.I.P.~System обслуживает онлайн-авторизацию с~жёстким таймаутом,
BASE~II принимает пакетные клиринговые инструкции в~собственном формате}.

Привязка \enquote{T+0~--- авторизация, T+1~--- клиринг} типична для рынка США
и~ряда крупных эквайеров, но это иллюстрация, а~не нормативная константа Visa
Core Rules: конкретные сроки клиринга и~расчёта зависят от~региона, типа
продукта и~участника.

Если эмитент недоступен или превышает таймаут ответа, Visa отвечает эквайеру
от~его имени через \textbf{Stand-In Processing (STIP)} на~основе параметров,
согласованных с~эмитентом (общий механизм~---
в~разделе~\ref{sec:tx-authorization}). Параллельно сетевой скоринг рисков
\textbf{Visa Advanced Authorization (VAA)} оценивает запрос с~помощью моделей
и~передаёт эмитенту оценку риска (см.~раздел~\ref{sec:visa-vaa}).

\begin{figure}[!ht]
  \centering
  \includefiguresvg[width=\linewidth]{assets/figures/ch13-visa/visanet-base-i-ii}
  \caption{Авторизационная сеть BASE~I и~расчётный бэкенд BASE~II.}\label{fig:visanet-base-i-ii}
\end{figure}

Для разработчика платёжного шлюза это две интеграции. Авторизационный запрос
передаётся в~реальном времени по~ISO~8583~--- международному стандарту
сообщений карточных авторизаций (см.~гл.~\ref{ch:iso8583}). Клиринговые записи
отправляются отдельным процессом по~расписанию; у~авторизации и~клиринга разные
классы ошибок и~регламенты обработки.

\section{Расчётная гарантия сети}\label{sec:visa-settlement-guarantee}

В~интервале между авторизацией и~расчётом эмитент уже подтвердил обязательство
заплатить, эквайер уже передал товар, но межбанковский неттинг через VSS
пройдёт через несколько часов или~суток. Если эмитент за~это время становится
несостоятельным, эквайер остаётся без~покрытия.

Правила Visa дают гарантию на~этот интервал: сеть обязуется возместить
финансовому институту-клиенту потери от~того, что другой клиент не~исполнил
расчётные обязательства~\cite{visa-10k-settlement-guarantee-fy2023}.
В~годовом отчёте
по~Form~10-K Visa называет это возмещение потерь при~расчёте
(\textit{settlement indemnification}) и~раскрывает методику оценки экспозиции:
среднедневной объём расчётов за~квартал, умноженный на~оценочное число дней
до~расчёта плюс защитная маржа, плюс объём возвратных платежей за~четыре
скользящих месяца~\cite{visa-10k-settlement-guarantee-fy2023}.

Первый уровень покрытия~--- обеспечение, которое сеть собирает с~участников,
не~прошедших кредитную оценку группы Global Credit Settlement Risk: денежные
депозиты в~эскроу на~имя Visa, заложенные ценные бумаги (хранятся у~кастодиана,
сеть вправе продать их при~дефолте), безотзывные аккредитивы банков клиента
и~гарантии материнской структуры по~обязательствам
дочерней~\cite{visa-10k-settlement-guarantee-fy2023,visa-canada-pps-self-assessment-2025}.
Когда обеспечения не~хватает, покрытие даёт собственная ликвидность сети: Visa
Global Treasury держит достаточно ликвидных активов, чтобы рассчитаться в~тот
же день даже при~дефолте участника с~наибольшей нетто-задолженностью,
и~ежедневно тестирует это требование по~стандартам
BIS~\cite{visa-canada-pps-self-assessment-2025}. Дополнительную опору даёт
собственный капитал Visa Inc.: на~30~сентября 2024~года $12{,}0$~млрд USD
в~денежных средствах и~их эквивалентах и~$5{,}4$~млрд USD в~долговых ценных
бумагах для~продажи (Treasury и~бумаги агентств правительства
США)~\cite{visa-canada-pps-self-assessment-2025}.

VCR/VPSR прямо запрещают одностороннюю отмену авторизованных, но ещё
не~рассчитанных операций даже в~случае несостоятельности
участника~\cite{visa-canada-pps-self-assessment-2025}. С~октября 2023~года
VisaNet формально обозначен Bank of Canada как Prominent Payment System
(системно значимая платёжная система) и~подчинён 18~стандартам управления
рисками, включая Standard~7 (финальность расчёта) и~Standard~9 (правила
и~процедуры при~дефолте участника)~\cite{visa-canada-pps-self-assessment-2025}.
Дополнительно к~основному механизму с~2007--2008~годов действует соглашение
о~распределении убытков между Visa Inc. и~банками-членами Visa U.S.A.: оно
распределяет покрытие отдельных категорий убытков по~доле членства независимо
от~стандартного возмещения при~расчёте.

У~Mastercard конструкция та~же: правила сети гарантируют расчёт по~значительной
части транзакций между клиентами, совокупный расчётный риск (\textit{gross
settlement exposure}) на~31~марта 2024~года оценён в~$73{,}8$~млрд USD
при~необеспеченном риске $61{,}2$~млрд, а~участники, не~проходящие стандарты
управления рисками, вносят денежное обеспечение, аккредитивы или
гарантии~\cite{mc-10q-settlement-risk-q1-2024}.

Механизм проявился в~марте 2023~года при~банкротствах Silicon Valley Bank
и~Signature Bank. На~конференции Wolfe Research финансовый директор Visa Васант
Прабху и~директор по~продуктам Mastercard Крейг Восбург подтвердили, что расчёт
через оба банка все выходные после перехода обоих банков под управление
регуляторов шёл без~сбоев, и~согласились с~описанием своих сетей как
страховщиков последней инстанции по~карточным
платежам~\cite{paymentsdive-svb-signature-card-networks-2023}.

В~ПС~Мир механизм аналогичный по~функции и~публично оформлен:
Приложение~3 к~Правилам ПС~Мир~--- \enquote{Методика определения
размера гарантийного обеспечения Участника}~--- задаёт формулу расчёта
обязательного депозита, который участник вносит в~НСПК пропорционально объёму
своих расчётов~\cite{nspk-mir-guarantee-fee-method}.

Карточный платёж в~четырёхсторонней модели свободен от~риска
контрагента-эмитента: при~банкротстве эмитента в~интервале между авторизацией
и~расчётом деньги эквайеру придут из~сети~--- сначала из~обеспечения этого
эмитента, затем из~ликвидной позиции сети, затем из~её капитала.
В~\textit{push}-системах вроде СБП, Pix или FedNow конструкция другая: расчёт
мгновенный и~идёт в~деньгах центрального банка, поэтому расчётного интервала
и~связанного с~ним кредитного риска по~отдельной транзакции там нет; вместо
него возникают риск ликвидности (покрытие на~корсчёте требуется в~момент
отправки) и~операционный риск (при~сбое системы платёж отклоняется немедленно:
исправить нельзя)~--- см.~гл.~\ref{ch:sbp}
и~гл.~\ref{ch:instant-payments-world}.

\section{Core Rules}\label{sec:visa-core-rules}

При споре, сертификации или разборе инцидента отправная точка~--- \textbf{свод
правил Visa (Visa Core Rules and Visa Product and Service Rules)}; его
публичная редакция доступна без~членского доступа, но Visa прямо предупреждает,
что из~неё исключены проприетарные сведения, коммерчески чувствительные данные
и~часть значимых для~безопасности деталей~\cite{visa-core-rules}.

Публикация устроена в~два уровня: \textbf{Visa Core Rules} задают общие
требования к~участию в~сети, ответственности сторон и~базовым процессам,
а~\textbf{Visa Product and Service Rules} добавляют продуктовые и~региональные
условия. Технические форматы сообщений, процедуры сертификации и~детали
сетевой обработки вынесены в~отдельные руководства по~внедрению (\textit{implementation
guides}), на~которые правила ссылаются, но содержание которых
не~пересказывают~\cite{visa-core-rules}.

За~пределами публичной редакции остаются членские соглашения с~конкретными
банками, дополнительные технические материалы, чувствительные сетевые параметры
и~индивидуальные коммерческие условия~--- они дополняют общий свод и~участнику
доступны через закрытые каналы.

Публичные редакции правил выходят с~фиксированной датой вступления в~силу
(\textit{effective date}); в~последние годы это апрельская и~октябрьская
редакции, и~к~каждой прилагается сводка изменений Summary of Changes.
Оперативные обновления удобно отслеживать через Visa Merchant Business News
Digest, который публикует краткие сводки и~отсылает к~материалам Business
News~\cite{visa-core-rules,visa-business-news-digest}. В~споре стороны
ссылаются на~редакцию, действовавшую на~дату операции~\cite{visa-core-rules}.

Для инженера Core Rules~--- рабочий документ участника: он показывает
обязанности эквайера и~эмитента, допустимые сроки клиринга, основания
для~понижения категории interchange и~момент, когда спор или инцидент
становятся нарушением правил сети
(см.~гл.~\ref{ch:regulation}).

\practicenote{%
  Читать Core Rules лучше от~задачи: чтение сверху вниз быстро
  превращается в~поверхностное знакомство. Сначала определите
  роль~--- эквайер, эмитент, VisaNet Processor (участник с~прямым
    подключением к~VisaNet, обрабатывающий трафик от/для~других членов
  сети) или конкретный продукт. Затем зафиксируйте редакцию документа
  и~дату вступления в~силу. После этого ищите локальную норму.
  Без~этой привязки легко сослаться на~правило, которое либо относится
  к~другой роли, либо ещё не~вступило в~силу в~нужном регионе.

}

\section{TADG и~сертификация терминалов}\label{sec:visa-tadg}

\textbf{Visa TADG (Terminal Acceptance Device Guide)}~--- \emph{инженерное
руководство Visa} по~поведению терминала, пользовательскому интерфейсу
и~сетевым деталям приёма карт; оно дополняет формальные спецификации и~тестовые
наборы, но не~заменяет их. Фраза \enquote{терминал сертифицирован под~Visa}
покрывает три проверки: аппаратную часть, ядро EMV и~интеграцию устройства
с~конкретным эквайером.

\textbf{Уровень 1 (L1)}~--- физический и~электрический уровень. Терминал
соответствует требованиям к~контактному считывателю, бесконтактному интерфейсу
и, при~наличии, магнитной полосе. Проверяются аппаратная часть и~низкоуровневое
взаимодействие с~картой~\cite{visa-tadg}.

\textbf{Уровень 2 (L2)}~--- программный уровень, ядро EMV. Терминал правильно
реализует логику транзакции EMV: выбор приложения, чтение данных, обработка
\textbf{TVR (Terminal Verification Results)}~--- вектора результатов проверок
терминала (см.~гл.~\ref{ch:emv}), генерация криптограмм. Для~контактных
устройств нужно одобрение консорциума EMVCo; для~части
бесконтактных стеков~--- соответствие \textbf{спецификации бесконтактных
платежей Visa (Visa Contactless Payment Specification, VCPS)}. На~L2 проверяют
реализацию сценария приёма целиком~\cite{visa-tadg}.

\textbf{Уровень 3 (L3)}~--- интеграционный уровень. Связка ядра EMV
с~конкретной системой управления терминалом (\textbf{Terminal Management
System, TMS}) эквайера и~его хостом авторизации. Устройство тестируют в~полевой
конфигурации: конкретное терминальное приложение, параметры эквайера,
маршрутизация авторизаций. На~L3 проверяют, как терминал формирует
авторизационное сообщение и~передаёт поле данных \textbf{DE~55} (Integrated
  Circuit Card data, контейнер EMV-данных в~сообщении ISO~8583,
см.~гл.~\ref{ch:iso8583}) с~данными EMV~\cite{visa-tadg}.

Отсюда практическое правило: пока не~назван уровень, слово
\enquote{сертифицирован} недостаточно. Устройство, не~поддерживающее нужный
сценарий EMV или введённое в~сеть с~неверной конфигурацией, лишает эквайера
защиты в~спорах о~фроде по~тому же механизму переноса ответственности
(\textit{liability shift}~--- перераспределение ответственности за~фрод между
эмитентом и~эквайером в~зависимости от~технологии приёма), что
и~в~разделе~\ref{sec:emv-magstripe-to-chip}.

\section{Visa Token Service}\label{sec:visa-vts}

Общая модель сетевых токенов EMVCo и~устройство сервиса \textbf{Visa Token
Service (VTS)} описаны в~главе~\ref{ch:tokenization}. Здесь~--- специфика Visa.

VTS предоставляет эмитенту три группы интерфейсов: \textbf{ID\&V}
(\textit{Identification \& Verification}~--- процедуры подтверждения личности
держателя при~выпуске токена), \textbf{Token
Inquiry} (запрос состояния токена и~данных об~устройстве и~риске)
и~\textbf{Token Lifecycle} (управление жизненным циклом токена) с~операциями
activate, resume, suspend и~delete~\cite{visa-vts}. Помимо интерфейсов Visa
вводит \textbf{Token Assurance Level} (уровень доверия к~токену): он
присваивается при~выпуске по~методу ID\&V (от~passive~--- без~активной
проверки, до~biometric~--- с~биометрическим подтверждением) и~передаётся
эмитенту в~каждой авторизации, влияя на~решение о~риске наряду с~\textbf{Token
Domain Restriction Controls}~--- ограничениями по~каналу, устройству и~типу
операции, для~которых выпущенный токен пригоден
(см.~раздел~\ref{sec:token-anatomy})~\cite{visa-vts,emvco-payment-tokenisation-guide}.

\section{Visa Direct}\label{sec:visa-direct}

Обычная карточная покупка движется по~модели \enquote{дебет у~плательщика},
в~которой авторизация резервирует средства, а~клиринг завершает списание. Visa
поддерживает и~обратный сценарий, когда деньги нужно доставить \emph{держателю
карты}. В~Visa Direct для~этого используется \textbf{исходная кредитовая
транзакция (Original Credit Transaction, OCT)}~--- зачисление на~счёт,
связанный с~картой получателя; \textbf{AFT (Account Funding Transaction)}
при~необходимости предварительно списывает средства с~карты
отправителя~\cite{visa-core-rules}.

OCT покрывает переводы между физлицами, массовые выплаты, пополнение
предоплаченных продуктов и~платежи по~кредитным картам~\cite{visa-direct}.

Принципиальное отличие OCT от~обычной авторизации~--- отсутствие холда
(резервирования суммы на~счёте плательщика до~завершения клиринга,
см.~гл.~\ref{ch:tx-lifecycle}): деньги кредитуются напрямую на~счёт получателя.
Публичная документация Visa оговаривает, что фактическая доступность средств
зависит от~банка получателя и~выбранного коридора. Общая норма Visa~---
средства зачисляются \enquote{как можно быстрее, но не~позднее двух рабочих
дней}; обещанный срок \enquote{до~30~минут} относится к~ускоренному сценарию
программы \textbf{Fast Funds} (добровольная программа Visa для~эмитентов
с~ускоренной доставкой средств получателю) и~применим только к~подключённым
к~ней эмитентам~\cite{visa-direct,visa-direct-faq}.

Типовой сценарий~--- связка OCT и~AFT. Если источник средств~--- карта
отправителя, инициатор выплаты сначала выполняет AFT, затем отправляет OCT
получателю. \highlight{AFT и~OCT~--- две отдельные транзакции: между ними
  может быть пауза в~минуты, часы или больше; в~это время платформа выплат
  может выполнять собственную проверку, ждать подтверждения источника или
менять источник финансирования}~\cite{visa-direct-faq}. Помимо связки
AFT$\rightarrow$OCT Visa допускает и~другие источники финансирования через Visa
Direct: дебет \textbf{ACH} (Automated Clearing House~--- американская сеть
межбанковских расчётов с~отложенным неттингом) или иной банковский перевод.
Такие сценарии не~оформляются как AFT и~подключаются по~отдельным программам.

\begin{figure}[!ht]
  \centering
  \includefiguresvg[width=\linewidth]{assets/figures/ch13-visa/visa-direct-flow}
  \caption{Типовой поток Visa Direct: AFT и~OCT.}\label{fig:visa-direct-flow}
\end{figure}

Клиринг по~AFT и~OCT передаётся отдельным потоком через BASE~II в~формате
TC~33.A~--- пакетном клиринговом файле Visa~\cite{visa-tc33a-capture-file};
схема на~рис.~\ref{fig:visa-direct-flow} отображает только авторизационный трек
V.I.P.~System.

Платформа выплат перед запуском фиксирует четыре параметра: кто финансирует
AFT, какой SLA (\textit{service-level agreement}, формальное соглашение
об~уровне обслуживания) установлен по~OCT, как банк получателя делает средства
доступными и~в~какой точке ответственность сети сменяется ответственностью
участника программы.

После марта 2022~года Visa Direct как трансграничный продукт для~российских
участников недоступен: эмиссия новых карт Visa в~санкционных банках прекращена,
а~приём и~отправка OCT/AFT за~пределы страны не~выполняются. Внутрироссийские
карточные операции по~картам Visa, выпущенным до~марта 2022~года, продолжают
обслуживаться через НСПК (см.~раздел~\ref{sec:mir-opkc}). Сбор средств между
своими счетами в~разных российских банках обслуживают Me2Me Pull и~Me2Me Push
(перевод самому себе с~инициацией из~принимающего и~отправляющего банка
соответственно), выплаты от~бизнеса физлицу~--- B2C-переводы
(\textit{business-to-consumer}) Системы быстрых платежей (СБП;
см.~гл.~\ref{ch:sbp} и~раздел~\ref{sec:sbp-scenarios}).

\section{Visa Advanced Authorization}\label{sec:visa-vaa}

Эмитент, принимающий решение по~подозрительной операции, видит только
собственный портфель: историю клиента, его поведение и~баланс. Весь сетевой
поток авторизаций доступен Visa, и~на~этом построен \textbf{сетевой скоринг
рисков Visa (Visa Advanced Authorization, VAA)}~\cite{visa-core-rules,visa-vaa}.

VAA оценивает авторизационный запрос до~решения эмитента: сеть рассчитывает
числовой показатель риска в~реальном времени и~передаёт его эмитенту в~составе
авторизационного сообщения~\cite{visa-vaa}.

\highlight{VAA работает на~уровне сети: Visa получает все
  авторизационные запросы от~эквайеров и,~прежде чем запрос дойдёт до~эмитента,
оценивает его с~помощью моделей, обученных на~сетевом трафике}. Модели
обнаруживают шаблоны, недоступные отдельному эмитенту,~--- например, что
по~той же карте только что прошли авторизации в~трёх разных городах.

В~одном варианте Visa Risk Manager принимает авторизационное решение автономно,
и~эмитент делегирует часть логики сети. В~другом сеть передаёт оценку риска
и~набор Risk Condition Codes (кодов выявленных рисковых условий, поясняющих, какие
сигналы повлияли на~оценку) в~авторизационный ответ, а~эмитент сам решает, как
реагировать,~--- снизить лимит сессии, запросить дополнительную аутентификацию
или одобрить операцию. Точное поле в~протоколе VisaNet для~этих данных
в~публичных материалах не~раскрыто; команда с~доступом к~VisaNet Message
Specifications сверяется по~закрытому документу. Сетевая оценка встраивается
в~собственную \textbf{RBA} эмитента (\textit{risk-based authentication}~---
  адаптивная аутентификация, при~которой требования к~подтверждению зависят
от~оценки риска операции) как ещё один сигнал наряду с~историей клиента
и~данными устройства.

\section{Interchange: структура таблиц}\label{sec:visa-interchange}

Общий механизм interchange, разложение MSC и~модели тарификации для~продавца
разобраны в~разделе~\ref{sec:ecosystem-interchange}; ниже~--- специфика Visa:
как организованы программы и~почему одна и~та~же операция может
квалифицироваться по~разным ставкам.

\highlight{Таблицы Visa сгруппированы по~\textbf{программам interchange},
  каждая программа~--- сочетание типа карты (дебетовая, кредитная,
  коммерческая), категории торгово-сервисного предприятия (ТСП) по~\textbf{MCC}
  (Merchant Category Code~--- четырёхзначный код категории ТСП по~ISO~18245)
  и~канала транзакции (карта предъявлена, электронная коммерция, бесконтактный
приём)}~\cite{visa-interchange}. ТСП с~MCC~5411 видит одну ставку для~дебетовых
карт, другую для~кредитных и~третью для~коммерческих; внутри этих веток заданы
дополнительные квалифицирующие условия по~данным и~способу приёма.

Механика \emph{downgrade}~--- общее свойство interchange, не~специфичное
для~Visa. В~Visa-таблицах оно
выражено через явные пары программ: транзакция, удовлетворяющая условиям лучшей
программы, квалифицируется как \enquote{qualified}; при~недостатке данных или
несоответствии сценария квалифицируется как \enquote{non-qualified} с~более
дорогой ставкой.

Таблица interchange Visa отражает интересы трёх сторон. Эмитенту выгодна
программа выше: каждая ступень дороже для~эквайера и~выгоднее эмитенту.
Эквайеру выгодна программа ниже: правильно собранный авторизационный
и~клиринговый пакет данных удерживает сделку в~дешёвой программе. Продавец
влияет только через MCC и~канал приёма, причём косвенно: MCC присваивает
эквайер по~Visa Merchant Data Standards
Manual~\cite{visa-merchant-data-standards-manual}. Региональные редакции таблиц
выходят дважды в~год, с~фиксированной датой вступления и~сводкой изменений;
ставка по~конкретной операции применяется на~дату обработки эквайером, а~не
на~дату покупки. Программа CPS в~США и~её аналоги в~других регионах~--- это
\enquote{лучшая} категория с~самыми жёсткими требованиями к~данным; переход
в~EIRF или Standard~IRF (Standard Interchange Reimbursement Fee~---
  максимальная по~стоимости категория downgrade Visa для~операций, не~прошедших
по~EIRF) означает, что одно из~условий не~выполнено.

Допустим, ТСП с~MCC~5912 (drugstore) принимает на~рынке США чиповую дебетовую
карту Visa на~\$100. При~корректной передаче DE~55 с~данными EMV транзакция
квалифицируется в~программе \textbf{CPS}\slash Retail (Custom Payment
  Service~--- группа программ interchange Visa с~расширенными требованиями
к~данным авторизации и~клиринга) со~ставкой порядка 0,80\,\% + \$0,15,
и~эквайер перечисляет эмитенту примерно \$0,95~\cite{visa-interchange}. Если
из-за ошибки терминала или хоста DE~55 не~попадает в~клиринговое сообщение,
операция квалифицируется по~категории \textbf{EIRF} (Electronic Interchange Reimbursement
  Fee~--- общая категория interchange Visa для~электронных операций
без~квалифицирующих признаков лучших программ) с~более высокой ставкой
около 1,75\,\% + \$0,20, и~interchange вырастает примерно до~\$1,95. Разница
порядка \$1,00 на~одной операции; при~10\,000~транзакций в~месяц downgrade
обходится продавцу около \$120\,000 в~год. Конкретные ставки и~потолки
сверяются по~актуальной редакции таблицы Visa Interchange Rates. Диагностика
downgrade начинается с~проверки, попал ли DE~55 в~клиринговый файл.
Авторизационный лог может показывать данные EMV, но если хост эквайера
не~передал их в~BASE~II, понижение всё равно
произойдёт~\cite{visa-core-rules,visa-tadg}.

Видимость down\-grade в~выписке зависит от~тарифной модели эквайера
(Inter\-change++, \textit{flat rate}, \textit{tiered}): в~Inter\-change++ понижение
программы видно сразу, в~\textit{flat rate} оно растворяется в~усреднённой
ставке.

\importantnote{%
  Потеря контекста EMV~--- это и~риск фрода, и~фактор interchange.
  Таблицы Visa различают \textit{card-present} (операция при~физическом
  предъявлении карты на~терминале), \textit{key-entry} (ручной ввод
  реквизитов оператором терминала) и~общие категории
  (EIRF, Standard~IRF); операция без~ожидаемых данных или вне
  сценария \textit{card-present} квалифицируется дороже. Конкретные
  ставки нужно сверять по~региональной таблице~\cite{visa-interchange,visa-tadg}.
}

\practicenote{%
  Чтобы найти программу \textit{interchange} для~конкретной транзакции
  в~документе Visa Interchange Rates, удобно двигаться от~общего
  к~частному. Сначала выбирают регион и~дату вступления в~силу таблицы.
  Затем определяют тип карты, категорию ТСП и~канал приёма. После
  этого читают квалифицирующие условия рядом со~ставкой~--- какие
  именно данные должны быть переданы в~сообщении, нужен ли контекст EMV,
  есть ли отдельные условия для~электронной коммерции или
  бесконтактного приёма.

}
